\documentclass[10pt]{article}
\usepackage[a4paper,margin=23mm]{geometry}
\usepackage{amsmath,amssymb,booktabs,graphicx,hyperref}
\setlength{\emergencystretch}{2em}
\title{Supplement: Controlled Response Experiments and Diagnostics}
\author{TMM working manuscript}
\date{October 18, 2026}
\begin{document}
\maketitle
These experiments explain the development and limits of the condensed method.
They are not extra modules attached to its inference pipeline. All results are
from the opened development cohort; complete protocols and negative results are
retained here rather than omitted during condensation.
\section{Directional post-refinement controls}
\paragraph{Directional post-refinement hypotheses.}
Let $r_0\in[0,1]^{N\times Q}$ denote initial element responses. Observation $e$
provides a fixed nonnegative row-normalized observation operator $A_e$, target
response $y_e$, bounded confidence $c_e\in[0,1]$, query compatibility $g_e$,
disjoint positive and exclusion permissions $p_e,n_e$, and an available
normalization domain $v_e\supseteq p_e\cup n_e$. Providers declare the provenance and allowed input
of these quantities. Scores are mapped to bounded responses by a frozen
provider; raw similarities from different teachers are not directly averaged.

Define $d_e=y_e-A_e r_0$ and
\begin{align}
 a_e^+ &= c_e g_e p_e\,\mathbf{1}[d_e>0], &
 a_e^- &= c_e g_e n_e\,\mathbf{1}[d_e<0],\\
 m &= \sum_e A_e^\top(a_e^++a_e^-), &
 b &= \sum_e A_e^\top[(a_e^++a_e^-)d_e].
\end{align}
The second post-refinement candidate uses the available geometric capacity
$D=\sum_e A_e^\top(g_e v_e)$ and performs one update:
\begin{equation}
 r^*_i=\begin{cases}
 \operatorname{clip}_{[0,1]}(r_{0,i}+b_i/\max(D_i,\epsilon)),&m_i>\epsilon,\\
 r_{0,i},&m_i\leq\epsilon,
 \end{cases}\qquad \epsilon=10^{-8}.
\end{equation}
All residuals use the original response; there is no iterative reuse of the
current prediction as new evidence. With nonnegative correspondences, positive
evidence alone cannot reduce the response. Explicit exclusion can reduce a
high initial response. Compatible unknown and zero-residual observations may
contribute normalization capacity, but unknown observations contribute no
gradient and rejected observations contribute neither. The available domain is
declared by the adapter; it is not every unobserved pixel in the scene.
These are algebraic properties, not guarantees of correctness of admitted
evidence or nondecreasing benchmark accuracy. Confidence scales the correction
without cancelling in its denominator; it is not a calibrated probability.
The first prototype used $D=m$. With one positive observation and nine unknown
observations of a scalar response 0.2, it can raise the response to 1.0;
the available-domain rule raises it only to 0.28. Both rules leave unknown
observations without an exclusion gradient. This is a normalization control
within the same admission mechanism, not a new feature representation.

The scientific focus is admission, not the weighted residual formula.
Query compatibility uses fixed native matching or explicit prompt support.
Missing detections, occlusion and rejected identities do not automatically
permit exclusion. A full reference mask and a partial detected region have
different exclusion authority. An inferred image anchor is not a hard 3D click.
The adapters now implement explicit permissions: LERF detections supply partial
positives; NVOS additionally permits exclusion from the retained reference-mask
estimate in its reference frame; the complete ScanNet category menu permits
exclusion from an unambiguous rival foreground. The latter two are fallible
provider estimates, not ground-truth negative prompts. UQIS uses frozen map
observations and identity only, averaging nested masks once per view. An anchor
inside a region is currently sufficient for positive support, an assumption
challenged by the diagnosis below.

For functional preservation, the implementation also exposes the retained
visibility fusion rule
\[
 r_{\mathrm{fusion}} = \frac{\sum_e A_e^\top(w_e y_e)}{\sum_e A_e^\top w_e},
\]
falling back to $r_0$ where there is no support. Its existing LERF entry point now calls
this extracted core. This compatibility path retains the original admission
policy and is distinct from the directional prototype above.

\section{Matched response studies and diagnosis}
Each study predicts all 31 jobs before scoring. The first three declare frozen
v10 responses as experimental inputs; the fourth uses the original pre-fusion
providers for source-assisted tasks. These are stage interventions, not fresh
end-to-end system runs. The baseline metrics reproduce the retained results.
The first study changes exclusion permission alone: partial evidence versus
allowing the available exterior as negative, with identical operators, targets
and confidence. Both use $D=m$. After inspecting this complete development
feedback, the second study preserves partial permissions and changes to $D$.
This latter control jointly changes normalization coverage and prevents
confidence cancellation; it does not separately identify their effects.

\begin{table}[ht]
\centering\small
\resizebox{\linewidth}{!}{\begin{tabular}{lrrrr}
\toprule
Metric (\%) & Frozen v10 & Active partial & Complement & Observable \\
\midrule
LERF2D IoU & 70.40 & 57.62 & 69.78 & 70.93 \\
LERF2D AP & 87.62 & 80.10 & 87.15 & 87.39 \\
NVOS IoU & 86.51 & 86.36 & 86.51 & 86.49 \\
NVOS AP & 98.02 & 97.70 & 97.99 & 97.99 \\
ScanNet mIoU19 & 61.01 & 58.30 & 58.07 & 60.59 \\
ScanNet mIoU15 & 62.53 & 59.28 & 59.14 & 61.88 \\
ScanNet mIoU10 & 72.49 & 70.30 & 69.93 & 72.03 \\
UQIS text AP & 26.60 & 24.13 & 22.34 & 25.98 \\
UQIS image AP & 59.31 & 54.37 & 50.81 & 57.19 \\
UQIS point\_2d AP & 75.79 & 67.18 & 64.16 & 71.93 \\
UQIS point\_3d AP & 75.33 & 67.07 & 64.00 & 72.09 \\
\bottomrule
\end{tabular}
}
\caption{All-cohort response controls on opened development scenes. Active
partial and complement use active-mass normalization; observable uses the same
partial permissions with available-domain normalization. No per-task selection.}
\label{tab:response-controls}
\end{table}

Active partial refinement severely harms LERF: it corrects 159,951 query pixels
but harms 2,185,722. Observable normalization reduces these to 36,202 and 35,192
and improves scene-macro IoU in all four scenes. The gain over v10 is 0.53 points
with a descriptive paired scene-bootstrap interval $[0.18,0.89]$; AP falls by
0.23 points. NVOS IoU changes by $-0.02$ points and ScanNet 19-class mIoU by
$-0.42$. UQIS text/image/2D-point/3D-point AP still falls by
0.62/2.12/3.86/3.24 points. Thus the normalization failure is actionable, but
the complete operator is not an accepted replacement.

LERF selected-set projection IoU remains 66.96. UQIS actual selected-set IoUs
remain 16.39, 39.90, 49.74 and 49.05 because the independent set heads are frozen.
UQIS sets are freshly rescored; LERF projection retains its unchanged control
score. These rows provide no evidence of improved set prediction.

Post-scoring evaluator-only diagnosis examines all 268 UQIS queries without
changing predictions. Both point modalities have 65 supported anchors correctly
located on the target instance, yet AP declines for 48 and 46 of these queries.
Non-target vertices receive 35.89\% and 30.62\% of their summed positive response
increments. This sum is a score-distribution diagnostic, not an error rate.
Even 35 of 44 image queries with correct estimated anchors lose AP. These
observations rule out wrong anchor identity as the sole explanation: a region
containing the target anchor can still supply inaccurate positive extent.
They motivate testing evidence admission at region and element level while
preserving the effective complete-membership provider. Labels used for this
diagnosis must never become query-time permissions.

\section{Scope agreement and further validation}
\paragraph{Scope agreement with a matched update budget.}
The third study keeps the previous permissions and capacity normalization.
Let $S$ denote the frozen provider's selected extent, or the full-menu category
argmax. A positive observation receives weight
$\rho_e^+=\sum (g_ep_e)(A_eS)/\max(\sum g_ep_e,1)$, computed per query;
negative observations analogously use $g_en_e$ and $1-A_eS$. The weight applies to
the whole observation, so support may still extend outside $S$. The guide does
not create negative permissions. A control uniformly attenuates the original
positive and negative update fields to match their respective per-query,
pre-clipping L1 masses. This controls total update size, not the number of
updated elements; clipping may change the final mass.

\begin{table}[ht]
\centering\small
\resizebox{\linewidth}{!}{\begin{tabular}{lrrrr}
\toprule
Metric (\%) & Frozen v10 & Observable & Scope & Matched \\
\midrule
LERF2D IoU & 70.40 & 70.93 & 70.55 & 70.54 \\
LERF2D AP & 87.62 & 87.39 & 87.63 & 87.60 \\
NVOS IoU & 86.51 & 86.49 & 86.49 & 86.49 \\
NVOS AP & 98.02 & 97.99 & 98.01 & 98.01 \\
ScanNet mIoU19 & 61.01 & 60.59 & 61.11 & 61.03 \\
ScanNet mIoU15 & 62.53 & 61.88 & 62.55 & 62.46 \\
ScanNet mIoU10 & 72.49 & 72.03 & 72.16 & 72.10 \\
UQIS text AP & 26.60 & 25.98 & 26.10 & 26.59 \\
UQIS image AP & 59.31 & 57.19 & 57.50 & 58.37 \\
UQIS point\_2d AP & 75.79 & 71.93 & 73.34 & 73.93 \\
UQIS point\_3d AP & 75.33 & 72.09 & 73.03 & 73.91 \\
\bottomrule
\end{tabular}
}
\caption{Third full-cohort study: extent agreement versus uniform attenuation
with matched signed correction budgets. Independent selected-set heads remain
fixed. Observable is the previous candidate, not a per-task fallback.}
\label{tab:support-agreement}
\end{table}

Scope agreement exceeds its matched control by only 0.004 LERF IoU points and
0.082 ScanNet 19-class mIoU points. On UQIS, text, image, 2D-point and 3D-point AP
are lower by 0.50, 0.87, 0.59 and 0.88 points. The image and point descriptive
scene-bootstrap intervals exclude zero. Thus improved accuracy relative to the
previous candidate cannot be assigned to scope-aware allocation alone. We do
not promote this candidate or assemble task-specific winners.

\paragraph{Where the intervention belongs.}
An audit of all 31 retained producer records finds that LERF and NVOS responses
already integrate the same masks, ScanNet already integrates source moments with
both native priors, and UQIS complete membership already uses anchor co-observation
features. The post-refinement studies therefore test reuse of consumed evidence.
This establishes information dependence, not the fraction of loss caused by
double counting. The fourth study therefore moves to response generation at
fixed inputs, preserving the two ScanNet teachers and UQIS complete
membership. Replacing a strong provider with raw similarity would not preserve
the functional baseline.

\paragraph{Matched intervention inside response generation.}
A fourth, predeclared control replaces response reconstruction itself rather
than adding another update to its output. All admitted targets and weights
$w_e$, native-positive vetoes, confidence scaling, teachers and independent
selected-set heads remain fixed. The diagonal consensus is
\begin{equation}
 r_{\mathrm{diag}}=
 \frac{\sum_e A_e^\top(w_e y_e)+\lambda\bar r}
 {\sum_e A_e^\top w_e+\lambda},
\end{equation}
with the original unsupported response retained when source mass is zero.
The rendered controls instead fit
\begin{equation}
 J(r)=\sum_{e,p,q}w_{e,pq}\rho\!\left(y_{e,pq}-(A_er)_{pq}\right)
       +\tfrac12\sum_{i,q}\lambda_{iq}(r_{iq}-\bar r_{iq})^2.
\end{equation}
One uses $\rho(t)=t^2/2$; the other uses a Huber derivative clipped at $0.25$.
Both take 120 bounded, visibility-preconditioned updates with step $0.95$,
inherited from the NVOS solver, without tuning on new scores. The squared
comparison separates diagonal assignment from forward mixture modeling;
the Huber comparison separately changes robustness. ScanNet retains its
original $\lambda=0.1$ times visible mass and the mean of both native priors;
other source protocols have no explicit prior penalty. LERF initializes from
its pre-fusion native response and NVOS from $0.5$. These inherited contracts
are disclosed rather than presented as a new universal calibration rule.
The prior penalty is evaluated on source-supported elements; other elements
remain fixed. The original solver is replayed alongside each intervention before evaluation.
UQIS has no corresponding raster reconstruction stage, so all its complete
membership outputs are explicitly unchanged compatibility controls. It cannot
establish transfer of this solver intervention to point or image membership.

\begin{table}[ht]
\centering\small
\resizebox{\linewidth}{!}{\begin{tabular}{lrrrr}
\toprule
Task metric & Frozen v10 & Diagonal & Rendered squared & Rendered huber \\
\midrule
LERF2D IoU & 70.40 & 70.40 & 70.80 & 70.82 \\
LERF2D AP & 87.62 & 87.62 & 87.00 & 86.58 \\
LERF selected projection IoU (fixed) & 66.96 & 66.96 & 66.96 & 66.96 \\
NVOS IoU & 86.51 & 86.57 & 86.52 & 86.51 \\
NVOS AP & 98.02 & 98.11 & 98.02 & 98.02 \\
ScanNet mIoU19 & 61.01 & 61.01 & 60.28 & 60.73 \\
ScanNet mIoU15 & 62.53 & 62.53 & 61.67 & 62.15 \\
ScanNet mIoU10 & 72.49 & 72.49 & 71.21 & 71.55 \\
UQIS text AP & 26.60 & 26.60 & 26.60 & 26.60 \\
UQIS text actual IoU (fixed) & 16.39 & 16.39 & 16.39 & 16.39 \\
UQIS image AP & 59.31 & 59.31 & 59.31 & 59.31 \\
UQIS image actual IoU (fixed) & 39.90 & 39.90 & 39.90 & 39.90 \\
UQIS point\_2d AP & 75.79 & 75.79 & 75.79 & 75.79 \\
UQIS point\_2d actual IoU (fixed) & 49.74 & 49.74 & 49.74 & 49.74 \\
UQIS point\_3d AP & 75.33 & 75.33 & 75.33 & 75.33 \\
UQIS point\_3d actual IoU (fixed) & 49.05 & 49.05 & 49.05 & 49.05 \\
\bottomrule
\end{tabular}
}
\caption{Fourth complete development study (\%): identical admitted source
observations and frozen providers, three reconstruction rules. All UQIS outputs
are unchanged compatibility controls; all independent selected sets are fixed.
No task-wise solver selection is applied.}
\label{tab:stage-reconstruction}
\end{table}

\paragraph{Fitting accuracy versus query quality.}
All 31 prediction reports were sealed and checked before scoring
(Table~\ref{tab:stage-reconstruction}). Squared reconstruction reduces the source
objective relative to diagonal consensus for all 22 source-assisted jobs.
Median relative reductions are 19.64\% for LERF, 20.94\% for NVOS, and 1.50\%
for ScanNet, whose objective includes its unchanged native-prior penalty.
LERF IoU increases in all four scenes by a mean 0.399 percentage points
(descriptive paired interval $[0.307,0.491]$), but AP decreases in all four
by a mean 0.616 points. NVOS changes are small: relative to diagonal consensus,
squared reconstruction changes IoU by $-0.057$ points, with interval
$[-0.198,0.105]$. ScanNet 19-class mIoU decreases by 0.732 points
($[-1.052,-0.386]$), with eight of ten rooms degrading. All ten rooms lose
10-class mIoU. Thus lower source objective is insufficient for improved target
quality even when the effective providers are preserved.

\begin{figure}[ht]
\centering
\includegraphics[width=\linewidth]{figures/stage_source_fit_vs_task.pdf}
\caption{Per-job source objective reduction and task score change, both relative
to diagonal consensus. Every source objective improves, but downstream changes
differ. Panels have separate horizontal scales and a common vertical scale;
the two NVOS horns requests are plotted separately. These are development data,
not independent validation.}
\label{fig:source-fit}
\end{figure}

Huber reconstruction recovers 0.446 ScanNet 19-class points over squared loss,
improving all ten rooms ($[0.302,0.601]$), but remains 0.286 points below
diagonal consensus. Its LERF IoU advantage over squared loss is only 0.023
points, with an interval spanning zero, while AP falls by another 0.422 points.
These results do not support declaring one more elaborate solver universally
better. An additive decomposition of the ScanNet scene-macro squared-loss
change finds negative contributions from 15 of the 18 classes present in
these rooms. Door, cabinet and wall contribute $-0.104$, $-0.101$ and $-0.074$
macro points, respectively. These are contributions to the macro change,
not the individual class IoU drops. They localize affected outputs but cannot
identify whether source identity, partial extent, geometry or view inconsistency
caused them. The next admission control must distinguish these hypotheses
without choosing observations using evaluation labels.

\paragraph{Further validation.}
The evaluation unit is the complete five-benchmark cohort. Each candidate
uses a fixed rule across scenes, preserves declared inputs, and seals its
predictions before aggregate scoring. The primary comparisons hold the teacher,
map and available evidence fixed: original response, existing simple fusion,
and the common refinement. Mechanism controls remove query admission, treat
unknown evidence as negative, restrict to positive evidence, or use matched
random rejection. Report final IoU/AP, harmful corrections, successful
corrections and unsupported-response abstention together.

Transfer experiments use at least two appropriate teacher configurations and
text, reference-image and point inputs. A teacher is only tested for modalities
it actually supports. UQIS 2D clicks cannot acquire source RGB, depth or masks
when attaching the refiner. A task that simply returns its prior is a valid
compatibility result, not evidence of improvement. New scenes must be held
apart from rule design; scene-level uncertainty and query latency/memory are
reported. Representation redesign and combinatorial system ablations are not
required for this response-level question.

\paragraph{Implementation status.}
The common core, protocol adapters and compatibility paths are implemented;
Targeted tests cover the numerical contracts. The production LERF fusion and NVOS robust solver call
extracted shared cores; ScanNet exposes its actual dual-teacher observations
through an opt-in capture interface. Default scientific rules are preserved.
In the fourth study, replaying each original solver gives maximum response
differences of $4.77\times10^{-7}$ for LERF, $4.72\times10^{-5}$ for NVOS and
$7.75\times10^{-7}$ for ScanNet. NVOS mean error never exceeds
$6.78\times10^{-8}$. There are no element-level threshold crossings at 0.5;
the original LERF/NVOS threshold predictions and ScanNet category predictions
are unchanged. Tiny AP drift from floating-point ranking is reported rather
than claimed as an improvement. A separate source-only audit compares the
literal old and extracted NVOS solvers within $1.08\times10^{-6}$ and confirms
318 source targets exactly match retained targets. The historical raster
operators were not retained, so the remaining historical discrepancy is not
uniquely attributed to CUDA accumulation.

Four complete studies and their paired feedback are available. The latest
study retains 1.98 GB of dense predictions only in volatile memory, releases
them after complete scoring, and adds approximately 5.6 MB of persistent
reports including failed, unscored attempts. The cumulative study remains
within its 2 GiB disk budget. Code snapshots, input identities, output hashes
and all results are retained; reproducing discarded arrays requires inference
again. Hashes bind the arrays scored in the run, not bitwise determinism across
GPU executions. Cross-teacher matched tests, random admission controls, improved
actual sets and an untouched validation cohort remain open. The retained v10
route remains the complete performance reference; the solver comparison is
an attribution experiment, not the final shared mechanism.

\section{Additional frozen-provider evidence}
\begin{table}[ht]
\centering\small
\resizebox{\linewidth}{!}{\begin{tabular}{lrrr}
\toprule
In-house component contrast & Before & After & $\Delta$ (pp) \\
\midrule
Dual-native category probabilities & 55.26 & 57.61 & +2.35 \\
Real-observation membership point\_2d & 40.12 & 49.74 & +9.62 \\
Real-observation membership point\_3d & 39.43 & 49.05 & +9.62 \\
\bottomrule
\end{tabular}
}
\caption{Additional evidence for frozen providers. The ScanNet comparison uses
three rooms held out during that historical study. Point membership uses the
same learner with different observation distributions and counts, and inherits
labelled-room supervision. These are not gains of the new shared operator.}
\end{table}

The point-membership contrast improves both point modalities in all nine
rooms. Its scene-bootstrap intervals are $[6.88,12.59]$ and $[6.53,12.68]$
percentage points for 2D and 3D points respectively. Keeping this provider
prevents the response study from discarding already demonstrated complete-object
capability. Its supervision and observation changes remain explicit.

The highest historical LERF semantic IoU in this response control is 70.75
from inverse rendering, with AP 86.58. Simple fusion achieves 70.38/87.64;
a native-response floor yields 68.57/89.11. We retain all variants and do not
stitch their metric-wise maxima into one method. Existing NVOS scope filtering
does not beat matched random filtering. ScanNet gains from an additional native
teacher are input-capability gains, not evidence for the new common operator.

\section{Estimated-anchor authority: a sparse causal control}
The effective image provider reuses the actual-point classifier and explicitly
sets its estimated anchor to one. We remove only this assignment, keeping the
conditioning anchor, model, inputs and calibrated threshold unchanged. All 67
image patches are fixed and hashed before evaluator access; all 268 UQIS queries
are then scored, including separate actual selections. No dense predictions are
saved. The remaining benchmark metrics are retained compatibility references,
not new inference or evidence of common-rule transfer.

Raw probabilities at the 67 anchors range from 0.91403 to 0.93902, all above
the frozen threshold 0.35. Exactly 67 score entries and zero selected elements
change. Image AP moves from 59.308782\% to 59.308926\%; the paired nine-scene
effect is +0.000144 points, with descriptive 95\% interval
$[-0.000759,0.001192]$. Actual image IoU remains 39.897960\%; text and actual
point scores and sets are exact. The override is therefore not a useful repair
target by itself on this cohort.

After prediction sealing, label-based diagnosis finds 44 correct and 23 incorrect
estimated anchors. Their image-IoU query means are 55.05\% and 9.89\%; conditional
scene means are 55.09\% across nine rooms and 10.75\% across eight rooms, respectively.
Sixteen incorrect anchors come from regions containing target elements; seven
selected regions contain none. Four correctly anchored queries still have IoU
below 0.25. These conditional comparisons neither constitute causal effects nor
justify an oracle switch. They distinguish the single-vertex override from the
more consequential, still unresolved identity condition of a complete response.
The raw self-member score has descriptive identity-correctness AUROC 0.5316 on
these 67 queries; its means are 0.9253 for correct and 0.9243 for incorrect anchors.
This analysis fits no threshold and supplies no query-time identity labels.
It indicates that this inherited confidence should not simply be reused as the
missing query-applicability estimate.

\section{Query witnesses before complete membership}
Q01 tests a query-conditioned identity interface while preserving the full-teacher
seed region and the complete classifier. It excludes the seed view, applies the
existing native-density identity gate to supported regions in other views,
normalizes their admitted weights per view, and chooses an actual seed member
nearest the resulting weighted coordinate median. The complete classifier and
threshold remain fixed. A control permutes the same regional weights within
each view, preserving the source-weight multiset and budget. This is one fixed
permutation, not a random-control seed sweep. The extra witness uses native
SigLIP cosine; original joint SigLIP/local-DINO seed selection is retained.
The candidate population differs from detector-conditioned LERF proposals.

All 67 original image responses replay exactly. Both intervention arms are
sealed before scoring all 268 queries. Other modalities remain exact; the
remaining benchmarks are historical references, not fresh all-31 inference.
\begin{table}[ht]
\centering\small
\begin{tabular}{lrrr}
\toprule
UQIS image output & v10 & Query witness & Permuted witness \\
\midrule
AP (\%) & 59.3088 & 58.7739 & 58.4641 \\
Actual-set IoU (\%) & 39.8980 & 38.7674 & 38.7483 \\
Correct anchors / 67 & 44 & 44 & 43 \\
\bottomrule
\end{tabular}

\caption{Q01: targeted identity-conditioning control on the opened development
cohort. No arm is promoted. Actual selections are re-evaluated.}
\end{table}

The query-conditioned arm changes 55 anchors and loses 1.1306 image-IoU points,
with descriptive nine-scene interval $[-2.4180,-0.1772]$; seven scenes decline.
Its advantage over the permutation control is only 0.0191 points, interval
$[-0.3648,0.4363]$. AP changes by $-0.5349$ points relative to v10, interval
$[-2.1267,0.5196]$. These results do not demonstrate query-specific benefit.
After sealing, one wrong anchor becomes correct and one correct anchor becomes
wrong. Among the 43 queries whose anchor stays on the target, query-mean IoU
falls from 55.8545\% to 54.3421\%. This conditional diagnosis is not a causal
comparison; it shows why correct-anchor count alone cannot verify preservation
of the complete response. The classifier is also sensitive to position and
observation evidence within the same instance. We retire this weighted-center
construction as a default candidate without threshold searches. Q01 retains
2.12 MB of reports and releases 117.81 MB of volatile prediction arrays.

\section{Verification among preserved complete responses}
Q02 keeps the original joint-teacher image answer and adds the SigLIP and local
DINO identity winners, each decoded by the same frozen complete classifier.
Cached equal-rank fusion is inverted on its checked half-integer rank lattice;
no DINO cosine magnitudes are reconstructed. A verifier compares the candidates'
co-visible regional IoU with query percentile evidence. In each common view,
it subtracts the IoU-complement-weighted query mean from the IoU-weighted mean.
The union of all proposal views is excluded. All candidates use the same
remaining informative views; no common evidence or a tie keeps the original.
One negative control permutes the same query percentiles within each view.
Both rules select one whole response without modifying it.

All 67 original image outputs replay exactly, and all proposals and choices
seal before scoring 268 queries. Other modalities remain exact; other benchmark
values are retained references. This is another targeted study, not a fifth
all-31 shared-method trial.
\begin{table}[ht]
\centering\small
\begin{tabular}{lrrr}
\toprule
UQIS image output & v10 & Verified & Permuted \\
\midrule
AP (\%) & 59.3088 & 59.5690 & 56.1093 \\
Actual-set IoU (\%) & 39.8980 & 38.2163 & 36.4864 \\
Correct anchors / 67 & 44 & 42 & 40 \\
\bottomrule
\end{tabular}

\caption{Q02: preserving complete proposals does not by itself guarantee a
correct identity-verification decision. Neither automatic rule is promoted.}
\end{table}

Image AP rises by 0.2602 points while actual IoU falls by 1.6816 points; the
latter descriptive nine-scene interval is $[-5.2721,2.1612]$. Relative to the
permutation control, IoU is higher by 1.7300 points, interval
$[-3.9031,6.6658]$, which does not establish a reliable query-specific effect.
The verifier repairs three wrong anchors but damages five formerly correct ones.
It retains the joint answer 30 times and chooses the other teachers 37 times.

After sealing, labels select a finite-pool best IoU of 45.0975\%; 27 queries
have a better proposal. This is a diagnostic, never an automatic score or an
upper bound on other proposal families. The exact scene-macro identity is
\begin{equation}
 I_{\rm automatic}=I_{\rm baseline}+G_{\rm pool}-R_{\rm selection},
\end{equation}
with 5.1995 points of available gain and 6.8811 points of selection regret.
Choosing the SigLIP or DINO winner globally gives only 33.5887\% and 31.1768\%.
Thus the finite candidates have useful diversity but the current verifier
cannot exploit it reliably. We retain the complete-proposal interface as
research infrastructure and reject its uncalibrated selector as a default.
Any calibrated replacement decision needs separate data and must measure both
repair and damage. Q02 stores 2.18 MB of reports and releases 294.52 MB of
volatile predictions; no new teacher weights or durable dense arrays are added.

\IfFileExists{figures/nvos_source_criteria.pdf}{
\section{Source-only NVOS decision calibration}
\begin{figure}[ht]
\centering\includegraphics[width=\linewidth]{figures/nvos_source_criteria.pdf}
\caption{N02 source criteria, not target-GT curves. Both estimators use the same
teacher targets and common pixels on every fourth source view. Fit-source
responses include those views in reconstruction; held-source responses exclude
their constraints. Dots denote the estimated thresholds and dashed vertical
lines denote 0.5. The final target response field is shared by all readout arms.}
\end{figure}
This criterion estimates a decision using source teacher evidence. Its maximum
does not certify target-view IoU or teacher correctness. Geometry and tracking
remain shared across the reconstruction split. The main manuscript reports the
complete eight-task results, including the fixed readout and both estimators.
}{}
\IfFileExists{../../output/tmm_response_20261016/study/nvos_score_changes_v1/errors.png}{
\section{Foreground and background score changes}
Let $O$ and $C$ be v10 and coverage responses. For diagnostic purposes only,
$P$ uses new scores on GT foreground and old scores on background; $N$ does the
opposite. We split the AP difference symmetrically:
\begin{align}
\Delta_{\rm fg}&=\tfrac12[\mathrm{AP}(P)-\mathrm{AP}(O)+\mathrm{AP}(C)-\mathrm{AP}(N)],\\
\Delta_{\rm bg}&=\tfrac12[\mathrm{AP}(N)-\mathrm{AP}(O)+\mathrm{AP}(C)-\mathrm{AP}(P)].
\end{align}
Their sum equals $\mathrm{AP}(C)-\mathrm{AP}(O)$ and divides the interaction
equally. These GT hybrids are not inference procedures. This is a descriptive
score-space attribution, not an upstream teacher or geometry causal ablation.
\begin{table}[ht]
\centering\small\begin{tabular}{lrrrr}
\toprule
Task & v10 AP & Coverage AP & Foreground $\Delta$ & Background $\Delta$ \\
\midrule
fern & 94.70 & 95.18 & +0.226 & +0.250 \\
flower & 99.32 & 99.39 & +0.090 & -0.027 \\
fortress & 99.98 & 99.99 & +0.003 & +0.007 \\
horns\_center & 99.78 & 99.81 & +0.056 & -0.022 \\
horns\_left & 98.87 & 99.01 & +0.100 & +0.042 \\
leaves & 99.79 & 99.86 & -0.022 & +0.090 \\
orchids & 97.83 & 95.81 & +0.221 & -2.237 \\
trex & 93.91 & 96.18 & +2.071 & +0.194 \\
\bottomrule
\end{tabular}

\caption{All eight tasks; AP in percent and score-group contributions in
percentage points. Actual response scores and GT are unchanged between audits.}
\end{table}
\clearpage
\section{Spatial changes on the complete NVOS cohort}
\begin{figure}[ht]
\centering\includegraphics[width=\linewidth,height=.79\textheight,keepaspectratio]{../../output/tmm_response_20261016/study/nvos_score_changes_v1/errors.png}
\caption{All eight NVOS tasks, using retained GT. The continuous field is freshly
rendered for v10; coverage and source-readout columns use the exact N02 response
images. Green denotes true positives, red false positives and blue false
negatives. The source-readout rule is fixed before target scoring. Display
resizing is for visualization only; metrics use the original GT resolution.}
\end{figure}
}{}
\IfFileExists{tables/nvos_local_credit.tex}{\section{Source-verified local response edits (N03)}
N02 tests a scalar decision scale; N03 instead asks where a response may change.
The anchor $b$ is our original v10 semantic response and $c$ is the fixed N01
coverage/high candidate. The same teachers and frozen geometry are reused.
Let $\delta=c-b$, $u=A\delta$, and $e=y-Ab$ on the common supported source domain.
For Huber loss $\rho$, define the integrated residual coefficient
\begin{equation}
 s_p=\begin{cases}
 [\rho(e_p)-\rho(e_p-u_p)]/u_p,&u_p\ne0,\\
 \rho'(e_p),&u_p=0.
 \end{cases}
\end{equation}
The signed and absolute allocations to element $i$ are
\begin{equation}
 g_i=\delta_i\sum_{v,p} A_{v,pi}w_{v,p}s_{v,p},\qquad
 h_i=|\delta_i|\sum_{v,p} A_{v,pi}w_{v,p}|s_{v,p}|.
\end{equation}
For nonnegative $A,w$, $|g_i|\le h_i$, and
$\sum_i g_i=\sum_{v,p}w_{v,p}[\rho(e_{v,p})-\rho(e_{v,p}-u_{v,p})]$.
This exact path accounting motivates a local gate
$q_i=\max(g_i,0)/h_i$, with zero permission for absent evidence. It is a
source-loss attribution identity, not a claim that the gated field has the same
loss reduction or higher target accuracy. A simple cancellation example shows why:
with $A=(1/2,1/2)$, $b=(0.5,0.5)$, $c=(1,0.1)$ and $y=0.55$, the candidate fits
exactly. Keeping only its positively credited first edit yields response 0.75,
which has higher source loss than the anchor's response 0.5. The quotient is an integrated loss
secant, not a newly claimed mathematical primitive.

Fit-source and held-source gates use exactly the same every-fourth eligible
source views, targets, weights and common support. The former uses the full
candidate; the latter uses a second field reconstructed without those views.
The latter gate applies to the final full-field edit only if its direction
agrees with the split-field edit; otherwise permission is zero. Geometry,
teacher tracks and the v10 anchor still use shared source information. Thus
this is a candidate-reconstruction holdout, not independent validation.

All anchored fields retain the v10 observation domain and output
$b+q\odot(c-b)$, preserving bounded values and the anchor where no update is
supported. A coverage-on-anchor-domain control isolates the domain change.
The matched uniform control separately preserves the positive and negative
Gaussian-space L1 edit masses for each query; it does not match projected
pixel counts. Six arms on all eight tasks yield 48 images sealed before target
scoring, with the unchanged 0.5 readout and both GT conventions. The held-local
arm is designated before scoring. No target RGB, GT-derived update regions,
per-task winners, or changed teachers are used.

\begin{table}[ht]
\centering\small
\resizebox{\linewidth}{!}{\begin{tabular}{lrrrrrr}
\toprule
Task & v10 & Coverage & Domain & Fit-local & Held-local & Uniform \\
\midrule
fern & 82.19 & 83.53 & 83.53 & 83.36 & 83.33 & 82.95 \\
flower & 86.53 & 87.38 & 87.38 & 87.26 & 87.25 & 86.95 \\
fortress & 97.23 & 97.46 & 97.46 & 97.42 & 97.42 & 97.53 \\
horns\_center & 95.73 & 95.55 & 95.55 & 95.89 & 95.90 & 95.85 \\
horns\_left & 92.70 & 94.43 & 94.43 & 94.16 & 94.11 & 94.00 \\
leaves & 94.08 & 95.06 & 95.06 & 95.00 & 95.00 & 95.00 \\
orchids & 68.14 & 70.05 & 70.05 & 69.64 & 69.61 & 69.04 \\
trex & 75.51 & 78.75 & 78.75 & 77.81 & 77.81 & 76.55 \\
\midrule Mean (retained GT) & 86.51 & 87.78 & 87.78 & 87.57 & 87.55 & 87.23 \\
\midrule Mean (LUDVIG GT) & 86.68 & 87.93 & 87.93 & 87.72 & 87.71 & 87.39 \\
\bottomrule
\end{tabular}
}
\caption{N03 all-eight IoU. Domain means coverage on the original anchor's
observed domain; uniform matches held-local signed edit mass.}
\label{tab:nvos-local-credit}
\end{table}
\begin{table}[ht]
\centering\small\begin{tabular}{lrrrr}
\toprule
Arm & IoU & Precision & Recall & AP \\
\midrule
v10 & 86.51 & 95.70 & 90.10 & 98.02 \\
Coverage & 87.78 & 95.83 & 91.43 & 98.15 \\
Anchor domain & 87.78 & 95.83 & 91.43 & 98.15 \\
Fit-local & 87.57 & 95.77 & 91.24 & 98.23 \\
Held-local & 87.55 & 95.76 & 91.23 & 98.27 \\
Matched uniform & 87.23 & 95.58 & 91.05 & 98.15 \\
\bottomrule
\end{tabular}

\caption{N03 macro metrics on retained GT. Every arm uses the same fixed
threshold; AP evaluates continuous scores.}
\label{tab:nvos-local-credit-macro}
\end{table}
The predeclared held-local arm obtains 87.5538\% IoU and 98.2704\% AP on retained GT.
Relative to the original v10, its mean IoU change is +1.0403 points (8 improved, 0 harmed), with a descriptive seven-scene bootstrap interval [+0.6636, +1.5292].
Relative to coverage, its mean IoU change is -0.2238 points (1 improved, 7 harmed), with a descriptive seven-scene bootstrap interval [-0.5000, -0.0479].
Relative to matched uniform attenuation, its mean IoU change is +0.3209 points (7 improved, 1 harmed), with a descriptive seven-scene bootstrap interval [+0.0750, +0.6771].
On the final fields, the common-source Huber objective is below the anchor for 8/8 held-local and 8/8 matched-uniform outputs. These fields use all source constraints, so this objective is not an independent validation score. The teacher-source criterion and target quality must be assessed separately.


The local allocation has a measured advantage over matched attenuation, but
it is not a replacement for the highest-IoU coverage candidate. Against
coverage, orchids recovers 1.4523 AP points while losing 0.4319 IoU points:
it gains 4,311 foreground pixels but loses 12,484, and adds 117 net false
positives. Trex loses 9,274 net foreground pixels and adds 9,129 net false
positives despite a lower mean background score. Relative to matched uniform
attenuation, however, the same trex local output gains 8,449 net foreground
pixels and removes 17,648 net false positives. The aggregate superiority over
uniform attenuation therefore does not establish preservation of every useful
coverage edit. All quantities are post-seal diagnostics, never update gates.

The observation-domain control changes mean IoU by only 0.000065 points.
Candidate-reconstruction holdout changes fit-local IoU by $-0.0141$ points;
we do not attribute an IoU gain to that extra reconstruction. Coverage remains
the principal best-IoU candidate, and there is no per-task selection between
coverage and local edits. Further constant or threshold searches around this
gate are not the next research step; the unresolved object extent and transport
of the admitted evidence require separate matched interventions.


\subsection{Does reference initialization discard the seed extent?}
A separate source-only audit compares the original predicted SAM reference mask
with the SAM3 track in that same authorized reference view. The fitted-query
lineage matches the declared query for every task. Agreement ranges from
96.49\% to 99.82\%; flower and orchids give 99.21\% and 97.47\% respectively.
This does not show a large reference-initialization collapse relative to the
seed. It does not establish that either mask describes the complete requested
object, or that later tracked views and geometric reconstruction are accurate.
Neither mask is ground truth; source agreement cannot substitute for target
accuracy or causally explain the target-RGB gap.
\begin{table}[ht]
\centering\small\begin{tabular}{lrrr}
\toprule
Task & Seed--track IoU & Seed retained & Track inside seed \\
\midrule
fern & 99.44 & 99.71 & 99.73 \\
flower & 99.21 & 99.57 & 99.64 \\
fortress & 99.82 & 99.92 & 99.90 \\
horns\_center & 99.56 & 99.72 & 99.83 \\
horns\_left & 99.67 & 99.83 & 99.84 \\
leaves & 99.81 & 99.88 & 99.93 \\
orchids & 97.47 & 98.37 & 99.07 \\
trex & 96.49 & 97.49 & 98.96 \\
\bottomrule
\end{tabular}

\caption{A22: agreement between two source predictions. ``Seed retained'' is
intersection/seed area, and ``track inside seed'' is intersection/track area.
These quantities are not benchmark recall or precision against ground truth.}
\end{table}
\begin{figure}[p]
\centering
\includegraphics[width=.94\linewidth,height=.87\textheight,keepaspectratio]{../../output/tmm_response_20261016/study/nvos_reference_scope_v1/reference_scope.jpg}
\caption{All eight authorized reference views, with original seed and tracked
extent. Difference panel: green overlap, blue seed-only, red track-only.
Target RGB and target GT are not used in this audit.}
\end{figure}
\IfFileExists{../../output/tmm_response_20261016/study/nvos_local_credit_spatial_v1/errors.png}{
\begin{figure}[p]
\centering
\includegraphics[width=\linewidth,height=.87\textheight,keepaspectratio]{../../output/tmm_response_20261016/study/nvos_local_credit_spatial_v1/errors.png}
\caption{N03 post-seal diagnostics on all eight tasks. Green: true positive;
red: false positive; blue: false negative. The held-local method and matched
uniform control use their predeclared rules, without task-wise selection.}
\end{figure}}{}
}{}
\IfFileExists{tables/nvos_image_response.tex}{\section{Image-conditioned query responses and 3D writeback}
N04 pairs a fixed v10 response with reconstructed RGB or observed RGB at the same camera. The former uses only source-trained appearance; the latter explicitly changes input access. SAM3 receives a response-derived box, optionally with dense log-odds. Single-mask decoding has no ground-truth candidate selection. Both dense-prior outputs are additionally written to the original observed Gaussian domain by 120 Huber updates and rendered again. Unsupported elements retain v10. The primary arm, observed RGB with dense prior and 3D writeback, was specified before scoring.
\begin{center}\small\begin{tabular}{lrrrrr}
\toprule
Arm & IoU & Precision & Recall & AP & LUDVIG GT IoU \\
\midrule
v10 & 86.51 & 95.70 & 90.10 & 98.02 & 86.68 \\
Render + box & 81.08 & 91.78 & 87.68 & 91.73 & 81.24 \\
RGB + box & 82.43 & 88.30 & 92.54 & 87.89 & 82.48 \\
Render + mask & 87.56 & 97.01 & 90.07 & 95.78 & 87.76 \\
RGB + mask & 89.06 & 96.13 & 92.51 & 95.04 & 89.13 \\
Render + mask + 3D & 86.92 & 97.05 & 89.35 & 96.78 & 87.16 \\
RGB + mask + 3D & 88.32 & 96.30 & 91.54 & 96.31 & 88.50 \\
\bottomrule
\end{tabular}
\end{center}
All 56 predictions precede target scoring. This is the opened development cohort, not an independent test. Neither target RGB nor the inferred image mask is target GT. A view-conditioned 3D update is not evidence of one query-global field or improved cross-view consistency. Source losses below compare against fallible original teacher tracks and are distinct from target accuracy.
The primary arm changes mean IoU by +1.8067 points versus v10 (6 improved, 2 harmed; descriptive seven-scene interval [+0.3180,+3.8393]).
The primary arm changes mean IoU by +1.3959 points versus the matched rendered-RGB arm (7 improved, 1 harmed; descriptive seven-scene interval [+0.6421,+2.3483]).
The primary arm changes mean IoU by -0.7435 points versus the direct image output (0 improved, 8 harmed; descriptive seven-scene interval [-1.0165,-0.5535]).
\begin{center}\small\begin{tabular}{lrrrrr}\toprule
Task & v10 IoU & RGB + 3D IoU & Source loss: v10 & Render + 3D & RGB + 3D \\\midrule
fern & 82.185 & 83.624 & 0.002636 & 0.009481 & 0.008375 \\
flower & 86.527 & 85.600 & 0.000511 & 0.003665 & 0.004451 \\
fortress & 97.230 & 96.985 & 0.000188 & 0.000728 & 0.000710 \\
horns\_center & 95.728 & 96.788 & 0.000441 & 0.002636 & 0.002938 \\
horns\_left & 92.698 & 93.212 & 0.000240 & 0.000451 & 0.000436 \\
leaves & 94.085 & 95.097 & 0.000202 & 0.000858 & 0.000974 \\
orchids & 68.141 & 72.669 & 0.000631 & 0.002021 & 0.003098 \\
trex & 75.513 & 82.586 & 0.002153 & 0.004970 & 0.004429 \\
\bottomrule\end{tabular}\end{center}
Implementation checks reproduce v10 binary outputs exactly, compare the sparse compositor and normalized response operator with gsplat, and bind RGB parameters to the exact geometry. Both GT conventions use continuous-score bilinear resizing before thresholding. An independent confusion-matrix implementation and scikit-learn AP check all arms. Other benchmark outputs are unchanged references, not fresh inference.
\IfFileExists{../../output/tmm_response_20261016/study/nvos_image_response_spatial_v2/errors.png}{
\begin{figure}[ht]\centering\includegraphics[width=\linewidth,height=.78\textheight,keepaspectratio]{../../output/tmm_response_20261016/study/nvos_image_response_spatial_v2/errors.png}
\caption{All eight N04 tasks. Green: true positive; red: false positive; blue: false negative. GT-conditioned diagnostics never enter inference.}\end{figure}}{}
}{}
\IfFileExists{tables/nvos_partial_response.tex}{\section{Positive identity and uncertain exclusion in image prompts}
N05 keeps the same response-derived box and eight deterministic inferred positive points for three image prompts: signed dense log-odds, positive-only dense log-odds (negative values become zero), and no dense prior. Point selection uses confidence, foreground interior distance and distance to previous points. These are estimated identity hints, not extra annotated clicks. The positive-only prompt with 3D writeback is the predeclared primary candidate. Neutral dense values are an operational SAM input, not a calibrated uncertainty estimate.
\begin{center}\small\begin{tabular}{lrrrrr}
\toprule
Arm & IoU & Precision & Recall & Accuracy & LUDVIG GT IoU \\\midrule
v10 & 86.51 & 95.70 & 90.10 & 97.55 & 86.68 \\
Signed + points & 88.72 & 95.88 & 92.38 & 97.90 & 88.79 \\
Positive + points & 85.93 & 97.11 & 88.33 & 97.55 & 86.09 \\
Points only & 84.07 & 89.66 & 93.97 & 96.67 & 84.07 \\
Signed + points + 3D & 88.02 & 96.07 & 91.43 & 97.80 & 88.19 \\
Positive + points + 3D & 85.38 & 97.26 & 87.64 & 97.47 & 85.61 \\
Points only + 3D & 83.66 & 90.02 & 93.04 & 96.63 & 83.77 \\
\bottomrule
\end{tabular}
\end{center}
All eight tasks use observed target RGB and the same frozen SAM3, original geometry, 512-by-384 lifting grid, 120 Huber steps and 0.5 threshold. Every prediction precedes scoring. N05 measures IoU, precision, recall and pixel accuracy under both GT conventions; AP is not measured. Results do not establish query-global multi-view consistency.
The primary candidate changes mean IoU by -1.1305 points relative to v10 (2 improved and 6 harmed tasks).
The primary candidate changes mean IoU by -2.6331 points relative to the signed dense control (0 improved and 8 harmed tasks).
The primary candidate changes mean IoU by +1.7232 points relative to the sparse-only control (3 improved and 5 harmed tasks).
\begin{center}\small\begin{tabular}{lrrrr}\toprule
Task & v10 & Signed + 3D & Positive + 3D & Sparse + 3D \\\midrule
fern & 82.19 & 83.78 & 83.45 & 77.20 \\
flower & 86.53 & 85.55 & 85.17 & 95.03 \\
fortress & 97.23 & 96.96 & 95.92 & 96.51 \\
horns\_center & 95.73 & 96.82 & 95.33 & 72.03 \\
horns\_left & 92.70 & 92.87 & 92.20 & 93.28 \\
leaves & 94.08 & 94.90 & 93.04 & 94.81 \\
orchids & 68.14 & 72.13 & 68.81 & 67.04 \\
trex & 75.51 & 81.11 & 69.13 & 73.38 \\
\bottomrule\end{tabular}\end{center}
\IfFileExists{../../output/tmm_response_20261016/study/nvos_partial_image_spatial_v1/errors.png}{
\begin{figure}[ht]\centering\includegraphics[width=\linewidth,height=.78\textheight,keepaspectratio]{../../output/tmm_response_20261016/study/nvos_partial_image_spatial_v1/errors.png}
\caption{All eight N05 tasks. Green: true positive; red: false positive; blue: false negative. No task-specific prompt switch is used.}\end{figure}}{}
}{}
\IfFileExists{tables/nvos_unboxed_response.tex}{\section{Query identity without inherited extent}
N06 keeps identical response-derived positive points and compares a box prompt, unboxed points, and unboxed points with signed dense log-odds. The primary candidate is unboxed sparse guidance with actual 3D writeback. It tests whether the old extent constrains completion. No new annotated clicks or target-dependent switch is used.
\begin{center}\small\begin{tabular}{lrrrrr}
\toprule
Arm & IoU & Precision & Recall & Accuracy & LUDVIG GT IoU \\\midrule
v10 & 86.51 & 95.70 & 90.10 & 97.55 & 86.68 \\
Box + points & 84.07 & 89.66 & 93.97 & 96.67 & 84.07 \\
Points, no box & 84.68 & 95.01 & 89.25 & 97.09 & 84.71 \\
Signed + points, no box & 87.93 & 96.87 & 90.58 & 97.82 & 88.03 \\
Box + points + 3D & 83.66 & 90.02 & 93.04 & 96.63 & 83.77 \\
Points, no box + 3D & 84.15 & 95.50 & 88.19 & 97.02 & 84.30 \\
Signed, no box + 3D & 87.16 & 97.00 & 89.66 & 97.70 & 87.36 \\
\bottomrule
\end{tabular}
\end{center}
All eight tasks use the same observed RGB, frozen SAM3, original Gaussian geometry, 120 Huber steps, 512-by-384 lifting grid and 0.5 threshold. All 56 predictions precede scoring. Both GT conventions are reported; AP is not measured. This is a development comparison and view-conditioned update, not a query-global field.
The primary candidate changes mean IoU by -2.3610 points relative to v10 (4 improved and 4 harmed tasks).
The primary candidate changes mean IoU by +0.4927 points relative to the boxed sparse control (3 improved and 5 harmed tasks).
The primary candidate changes mean IoU by -3.0088 points relative to the unboxed signed control (2 improved and 6 harmed tasks).
\begin{center}\small\begin{tabular}{lrrrr}\toprule
Task & v10 & Box + 3D & No box + 3D & Signed, no box + 3D \\\midrule
fern & 82.19 & 77.20 & 83.28 & 83.45 \\
flower & 86.53 & 95.03 & 72.12 & 85.41 \\
fortress & 97.23 & 96.51 & 96.01 & 96.76 \\
horns\_center & 95.73 & 72.03 & 95.25 & 96.53 \\
horns\_left & 92.70 & 93.28 & 93.00 & 92.75 \\
leaves & 94.08 & 94.81 & 94.77 & 94.78 \\
orchids & 68.14 & 67.04 & 58.65 & 71.24 \\
trex & 75.51 & 73.38 & 80.15 & 76.37 \\
\bottomrule\end{tabular}\end{center}
\IfFileExists{../../output/tmm_response_20261016/study/nvos_unboxed_image_spatial_v1/errors.png}{
\begin{figure}[ht]\centering\includegraphics[width=\linewidth,height=.78\textheight,keepaspectratio]{../../output/tmm_response_20261016/study/nvos_unboxed_image_spatial_v1/errors.png}
\caption{All eight N06 tasks. Green: true positive; red: false positive; blue: false negative. No task-specific prompt switch is used.}\end{figure}}{}
}{}
\IfFileExists{tables/nvos_memory_response.tex}{\section{Reference memory response transport}
N07 carries the original positive-only reference mask through a frozen SAM3 tracker. Two-frame memory uses the reference and observed target image; sequence memory also uses the ordered source views between them. The primary arm is sequence memory with actual 3D writeback. The source mask is fallible, and temporal memory is an existing external teacher capability, not a claimed new tracker.
\begin{center}\small\begin{tabular}{lrrrrr}
\toprule
Arm & IoU & Precision & Recall & Accuracy & LUDVIG GT IoU \\\midrule
v10 & 86.51 & 95.70 & 90.10 & 97.55 & 86.68 \\
Static image & 89.06 & 96.13 & 92.51 & 97.95 & 89.13 \\
Pair memory & 88.40 & 94.97 & 93.10 & 97.83 & 88.39 \\
Sequence memory & 88.35 & 95.04 & 92.99 & 97.81 & 88.35 \\
Static image + 3D & 88.32 & 96.30 & 91.54 & 97.84 & 88.50 \\
Pair memory + 3D & 87.94 & 95.37 & 92.16 & 97.77 & 88.08 \\
Sequence memory + 3D & 87.89 & 95.42 & 92.06 & 97.76 & 88.03 \\
\bottomrule
\end{tabular}
\end{center}
The static N04 image response is freshly replayed. Each arm uses the same image checkpoint, image preprocessing, original Gaussian geometry, 120 Huber updates, 512-by-384 lifting grid and 0.5 threshold. All 56 predictions precede both GT scorings. AP is not measured. No task-wise switching is allowed. These are development results with observed target RGB; the updated fields are view-conditioned, not query-global.
The primary candidate changes mean IoU by +1.3718 points relative to v10 (7 improved and 1 harmed tasks).
The primary candidate changes mean IoU by -0.4349 points relative to the static control (5 improved and 3 harmed tasks).
The primary candidate changes mean IoU by -0.0597 points relative to two-frame memory (3 improved and 4 harmed tasks).
\begin{center}\small\begin{tabular}{lrrrr}\toprule
Task & v10 & Static + 3D & Pair + 3D & Sequence + 3D \\\midrule
fern & 82.19 & 83.62 & 83.89 & 84.21 \\
flower & 86.53 & 85.60 & 88.75 & 87.51 \\
fortress & 97.23 & 96.99 & 97.32 & 97.21 \\
horns\_center & 95.73 & 96.79 & 95.18 & 97.38 \\
horns\_left & 92.70 & 93.21 & 94.52 & 94.55 \\
leaves & 94.08 & 95.10 & 95.26 & 95.01 \\
orchids & 68.14 & 72.67 & 70.02 & 70.02 \\
trex & 75.51 & 82.59 & 78.63 & 77.18 \\
\bottomrule\end{tabular}\end{center}
\IfFileExists{../../output/tmm_response_20261016/study/nvos_memory_spatial_v1/errors.png}{
\begin{figure}[ht]\centering\includegraphics[width=\linewidth,height=.78\textheight,keepaspectratio]{../../output/tmm_response_20261016/study/nvos_memory_spatial_v1/errors.png}
\caption{All eight N07 tasks. Green: true positive; red: false positive; blue: false negative. Original reference identity is transferred with the same rule across tasks.}\end{figure}}{}
}{}
\IfFileExists{tables/nvos_native_writeback.tex}{\section{Resolution-consistent response transport}
N08 keeps every image probability from N07 fixed and changes the writeback ray grid from 512-by-384 to the native 1008-by-756 output grid. The Gaussian geometry and observed domain, initial v10 field, absolute-margin confidence, 120 Huber steps and 0.5 threshold are unchanged. The three original coarse writebacks are exact frozen controls. All 56 responses precede target scoring; both GT conventions are reported. The primary arm is sequence-memory native writeback. This design was registered before any N07 scores. AP is not measured. Results remain view-conditioned development evidence.
\begin{center}\small\begin{tabular}{lrrrrr}
\toprule
Writeback & IoU & Precision & Recall & Accuracy & LUDVIG GT IoU \\\midrule
v10 & 86.51 & 95.70 & 90.10 & 97.55 & 86.68 \\
Static, coarse & 88.32 & 96.30 & 91.54 & 97.84 & 88.50 \\
Pair, coarse & 87.94 & 95.37 & 92.16 & 97.77 & 88.08 \\
Sequence, coarse & 87.89 & 95.42 & 92.06 & 97.76 & 88.03 \\
Static, native & 88.84 & 96.21 & 92.20 & 97.91 & 88.96 \\
Pair, native & 88.20 & 95.10 & 92.73 & 97.80 & 88.27 \\
Sequence, native & 88.11 & 95.13 & 92.62 & 97.78 & 88.19 \\
\bottomrule
\end{tabular}
\end{center}
static image: native writeback changes IoU by +0.5223 points (8 improved, 0 harmed tasks).
pair memory: native writeback changes IoU by +0.2553 points (7 improved, 1 harmed tasks).
sequence memory: native writeback changes IoU by +0.2272 points (7 improved, 1 harmed tasks).
\begin{center}\small\begin{tabular}{lrrrr}\toprule
Task & v10 & Static native & Pair native & Sequence native \\\midrule
fern & 82.19 & 83.92 & 84.06 & 84.33 \\
flower & 86.53 & 85.98 & 89.25 & 87.91 \\
fortress & 97.23 & 97.36 & 97.66 & 97.56 \\
horns\_center & 95.73 & 97.30 & 95.47 & 97.79 \\
horns\_left & 92.70 & 93.68 & 94.80 & 94.86 \\
leaves & 94.08 & 95.61 & 95.68 & 95.53 \\
orchids & 68.14 & 73.17 & 70.59 & 70.59 \\
trex & 75.51 & 83.71 & 78.09 & 76.33 \\
\bottomrule\end{tabular}\end{center}
\begin{figure}[ht]\centering\includegraphics[width=\linewidth,height=.78\textheight,keepaspectratio]{../../output/tmm_response_20261016/study/nvos_native_spatial_v2/errors.png}
\caption{N08 complete cohort: green = TP, red = FP, blue = FN. The sequence coarse column uses the exact N07 response.}\end{figure}
\paragraph{Readout capacity, not a method result.}
An exhaustive threshold curve refines the earlier A15 grid diagnosis. Equal scores are kept together. Even one GT-optimal threshold per task is a nondeployable upper bound on scalar-threshold decisions of unchanged v10 responses. It is not an upper bound on new spatially conditioned responses, and no oracle threshold is exported.
\begin{center}\small\begin{tabular}{lrrr}\toprule
Task & Fixed v10 IoU & GT-oracle IoU & AP \\\midrule
fern & 82.19 & 83.41 & 94.70 \\
flower & 86.53 & 91.34 & 99.32 \\
fortress & 97.23 & 98.47 & 99.98 \\
horns\_center & 95.73 & 96.21 & 99.78 \\
horns\_left & 92.70 & 93.05 & 98.87 \\
leaves & 94.08 & 96.00 & 99.79 \\
orchids & 68.14 & 85.83 & 97.83 \\
trex & 75.51 & 76.56 & 93.91 \\
Mean & 86.51 & 90.11 & 98.02 \\
\bottomrule\end{tabular}\end{center}
The optional CUDA footprint backend corrects eager arithmetic cancellation on thin conics; this is an implementation correction, not a new segmentation mechanism. The matched coarse-grid backend control changes cohort IoU by at most 0.0000053 percentage points. The full native run retains the original renderer-consistency tolerance.
}{}
\IfFileExists{tables/nvos_signed_affinity.tex}{\section{Signed reference evidence without an extent gate}
N09 uses frozen DINOv2 ViT-B/14 descriptors on the original reference and target RGB. The primary branch gives authority only to original positive and negative scribbles; unlabelled reference pixels are unknown. The matched pseudo-mask control adds the retained reference SAM mask/complement as labels and applies the same scribble overrides. Both branches use identical teachers, query inputs and downstream operations. This explicitly extends the positive-only v10 input contract, without feature-field training. DINOv2 and SAM are external teachers, not new contributions of this paper.
For normalized reference descriptors $z_j$, define $m_i=\max_{j\in P}z_i^\top z_j-\max_{j\in N}z_i^\top z_j$ and $p_i=\sigma(m_i/\tau)$. The scale $\tau$ is the median signed leave-one-out reference margin divided by $\log 9$. Mixed-sign patches are unknown. Missing support or a nonpositive median disables the edit and uses the static response. The frozen SAM decoder receives the transferred probability and its own extent box. Native-ray 3D writeback uses the same 120-step solver and observed domain as N08. No target-selected threshold, per-task winner or new user annotation is used.
All 72 responses across eight tasks precede GT. Both GT conventions and AP are scored. Other benchmarks are unchanged references; outputs remain view-conditioned development evidence.
\begin{center}\small\resizebox{\linewidth}{!}{\begin{tabular}{lrrrrr}
\toprule
Variant & IoU & AP & Precision & Recall & LUDVIG GT IoU \\\midrule
v10 & 86.51 & 98.02 & 95.70 & 90.10 & 86.68 \\
Static image & 89.06 & 95.04 & 96.13 & 92.51 & 89.13 \\
Static native & 88.84 & 96.58 & 96.21 & 92.20 & 88.96 \\
Signed affinity & 77.52 & 94.96 & 83.06 & 92.69 & 77.35 \\
Pseudo affinity & 84.10 & 96.72 & 93.43 & 89.49 & 84.11 \\
Signed + SAM & 81.68 & 86.27 & 85.07 & 95.15 & 81.62 \\
Pseudo + SAM & 86.23 & 93.87 & 90.85 & 94.73 & 86.22 \\
Signed native (primary) & 81.58 & 87.72 & 85.25 & 94.88 & 81.60 \\
Pseudo native & 86.08 & 95.94 & 90.99 & 94.42 & 86.14 \\
\bottomrule
\end{tabular}
}\end{center}
\begin{center}\small\begin{tabular}{lrrrr}\toprule
Task & Static native & Signed native & Pseudo native & Signed direct affinity \\\midrule
fern & 83.92 & 83.21 & 83.64 & 76.40 \\
flower & 85.98 & 95.07 & 94.89 & 89.00 \\
fortress & 97.36 & 97.88 & 97.62 & 88.00 \\
horns\_center & 97.30 & 67.16 & 95.80 & 73.76 \\
horns\_left & 93.68 & 91.07 & 93.90 & 72.44 \\
leaves & 95.61 & 45.57 & 66.91 & 70.53 \\
orchids & 73.17 & 88.89 & 76.04 & 83.08 \\
trex & 83.71 & 83.79 & 79.85 & 66.92 \\
\bottomrule\end{tabular}\end{center}
signed native vs static native: -7.2621 IoU points, -8.8662 AP points; 4 improved and 4 harmed tasks.
signed native vs pseudo native: -4.5009 IoU points, -8.2269 AP points; 4 improved and 4 harmed tasks.
signed affinity vs pseudo affinity: -6.5841 IoU points, -1.7692 AP points; 1 improved and 7 harmed tasks.
signed image vs signed affinity: +4.1607 IoU points, -8.6843 AP points; 6 improved and 2 harmed tasks.
pseudo image vs pseudo affinity: +2.1325 IoU points, -2.8508 AP points; 7 improved and 1 harmed tasks.
signed native vs signed image: -0.0975 IoU points, +1.4444 AP points; 2 improved and 6 harmed tasks.
pseudo native vs pseudo image: -0.1526 IoU points, +2.0685 AP points; 0 improved and 8 harmed tasks.
\begin{figure}[ht]\centering\includegraphics[width=\linewidth,height=.78\textheight,keepaspectratio]{../../output/tmm_response_20261016/study/nvos_signed_affinity_v1/errors.png}
\caption{N09 full cohort: green TP, red FP, blue FN.}\end{figure}
}{}
\IfFileExists{tables/nvos_direct_affinity.tex}{\section{Direct affinity writeback without a second segmentation}
N10 was fixed before N09 scoring. Exact signed and pseudo-mask affinity probabilities from N09 are fitted directly into the same observed Gaussians using the identical native-ray operator, confidence and 120-step objective. The four parent controls retain exact continuous-array hashes and scores. No teacher calls, new prompts, geometry or target-selected parameters are introduced. All 48 responses precede GT. Original negative reference scribbles and observed target RGB remain explicit input extensions; the result is view-conditioned development evidence.
\begin{center}\small\resizebox{\linewidth}{!}{\begin{tabular}{lrrrrr}
\toprule
Variant & IoU & AP & Precision & Recall & LUDVIG GT IoU \\\midrule
v10 & 86.51 & 98.02 & 95.70 & 90.10 & 86.68 \\
Static SAM native & 88.84 & 96.58 & 96.21 & 92.20 & 88.96 \\
Signed + SAM native & 81.58 & 87.72 & 85.25 & 94.88 & 81.60 \\
Pseudo + SAM native & 86.08 & 95.94 & 90.99 & 94.42 & 86.14 \\
Signed direct native (primary) & 77.73 & 95.08 & 83.34 & 92.65 & 77.58 \\
Pseudo direct native & 84.21 & 96.78 & 93.62 & 89.44 & 84.22 \\
\bottomrule
\end{tabular}
}\end{center}
\begin{center}\small\begin{tabular}{lrrrr}\toprule
Task & Static & Signed + SAM & Signed direct & Pseudo direct \\\midrule
fern & 83.92 & 83.21 & 76.40 & 80.64 \\
flower & 85.98 & 95.07 & 89.46 & 91.38 \\
fortress & 97.36 & 97.88 & 88.02 & 94.45 \\
horns\_center & 97.30 & 67.16 & 73.90 & 94.01 \\
horns\_left & 93.68 & 91.07 & 72.72 & 87.81 \\
leaves & 95.61 & 45.57 & 70.55 & 84.47 \\
orchids & 73.17 & 88.89 & 83.64 & 69.93 \\
trex & 83.71 & 83.79 & 67.15 & 70.99 \\
\bottomrule\end{tabular}\end{center}
signed direct native vs static native: -11.1117 IoU points, -1.5029 AP points; 2 improved and 6 harmed tasks.
signed direct native vs signed native: -3.8496 IoU points, +7.3633 AP points; 2 improved and 6 harmed tasks.
pseudo direct native vs pseudo native: -1.8713 IoU points, +0.8362 AP points; 1 improved and 7 harmed tasks.
signed direct native vs pseudo direct native: -6.4791 IoU points, -1.6998 AP points; 1 improved and 7 harmed tasks.
}{}
\IfFileExists{tables/nvos_instance_extent.tex}{\section{Instance-constrained additive extent recovery}
N11 retains the exact N09 static and signed-SAM image evidence. Three matched policies accept all proposal components, use the v10 response as signed identity evidence, or use original signed scribbles transported through fixed geometry (primary). Unmarked reference coverage has zero authority in the geometry branch. Eight-connected proposal regions require positive evidence mass greater than negative mass; empty/tied evidence retains the static response. Admitted pixels take the maximum of static and proposal probability. This image update is additive, but downstream 3D fitting is not guaranteed monotone. Every branch uses the same native rays and 120-step fit. DINOv2, SAM3, target RGB and official negative scribbles are explicitly retained inputs. All 72 predictions precede GT; other benchmarks remain unchanged references.
\begin{center}\small\resizebox{\linewidth}{!}{\begin{tabular}{lrrrrr}
\toprule
Variant & IoU & AP & Precision & Recall & LUDVIG GT IoU \\\midrule
v10 & 86.51 & 98.02 & 95.70 & 90.10 & 86.68 \\
Static image & 89.06 & 95.04 & 96.13 & 92.51 & 89.13 \\
Static native & 88.84 & 96.58 & 96.21 & 92.20 & 88.96 \\
Unrestricted image & 82.96 & 94.23 & 85.64 & 96.59 & 82.90 \\
Unrestricted native & 83.58 & 96.81 & 86.80 & 96.00 & 83.64 \\
v10-anchor image & 90.19 & 94.60 & 95.16 & 94.65 & 90.13 \\
v10-anchor native & 89.97 & 96.43 & 95.43 & 94.17 & 90.04 \\
Scribble-geometry image & 90.16 & 94.60 & 95.13 & 94.66 & 90.10 \\
Scribble-geometry native (primary) & 89.97 & 96.43 & 95.43 & 94.18 & 90.05 \\
\bottomrule
\end{tabular}
}\end{center}
\begin{center}\small\begin{tabular}{lrrrr}\toprule
Task & Static & Unrestricted & v10 anchor & Scribble geometry \\\midrule
fern & 83.92 & 83.73 & 83.67 & 83.66 \\
flower & 85.98 & 94.93 & 94.93 & 94.93 \\
fortress & 97.36 & 97.77 & 97.77 & 97.77 \\
horns\_center & 97.30 & 75.49 & 97.42 & 97.43 \\
horns\_left & 93.68 & 93.30 & 93.70 & 93.70 \\
leaves & 95.61 & 52.49 & 95.66 & 95.66 \\
orchids & 73.17 & 87.82 & 73.54 & 73.54 \\
trex & 83.71 & 83.13 & 83.05 & 83.10 \\
\bottomrule\end{tabular}\end{center}
scribble anchor native vs static native: +1.1310 IoU points, -0.1540 AP points; 6 improved and 2 harmed tasks.
scribble anchor native vs union native: +6.3929 IoU points, -0.3869 AP points; 3 improved and 3 harmed tasks.
scribble anchor native vs prior anchor native: +0.0061 IoU points, +0.0003 AP points; 4 improved and 1 harmed tasks.
prior anchor native vs static native: +1.1248 IoU points, -0.1544 AP points; 6 improved and 2 harmed tasks.
union native vs static native: -5.2619 IoU points, +0.2329 AP points; 3 improved and 5 harmed tasks.
scribble anchor native vs scribble anchor image: -0.1884 IoU points, +1.8326 AP points; 3 improved and 5 harmed tasks.
\begin{figure}[ht]\centering\includegraphics[width=.85\linewidth,height=.75\textheight,keepaspectratio]{../../output/tmm_response_20261016/study/nvos_instance_extent_spatial_v1/errors.png}
\caption{All eight N11 tasks. Green TP, red FP, blue FN; native-grid display resizing only. Scores use original GT resolution.}\end{figure}
\paragraph{Instance-constrained extent proposals.}
The N11 extension preserves the static image response $p_t$ while admitting
additional regions from a fixed signed-reference proposal $q_t$. For reference
view $r$ and raw visible Gaussian contributions $C_r$, original signed scribbles
$P,N$ give $b^+=C_r^\top P/(C_r^\top\mathbf1)$ and
$b^-=C_r^\top N/(C_r^\top\mathbf1)$, with zero where the denominator is empty.
Unmarked reference coverage supplies neither sign. Restricting to the original
observed domain, render $W_t^\pm=\bar A_t b^\pm$. For each eight-connected
component $M_k$ of $q_t\geq0.5$, define
\[
 a_k=\mathbf1\!\left[\sum_{u\in M_k}W_t^+(u)>
 \sum_{u\in M_k}W_t^-(u),\quad \sum_{u\in M_k}W_t^+(u)>10^{-8}\right].
\]
Accepted pixels take $\max(p_t,q_t)$; all other pixels retain $p_t$. This
changes image support authority before the same native-ray fit; it guarantees
no image-score decrease, not monotonic 3D predictions. Unrestricted admission
and $W_t^+=z_t^0,W_t^-=1-z_t^0$ are matched controls. Original negative scribbles
and frozen DINOv2 descriptors are declared additions to the v10 input contract.
The experiments support instance constraints, but do not establish an advantage
of this geometric witness over the retained v10 response anchor.

}{}
\IfFileExists{tables/nvos_reference_completion.tex}{\section{Signed reference completion before instance transport}
N12 compares the original reference mask, a frozen-DINO signed-affinity response decoded on the reference image, and the same completion with original positive/negative scribble points plus exact annotation constraints. All three use identical frozen two-frame SAM3 memory and native-ray 3D fitting. The primary constrained branch is fixed for all eight tasks; all 72 predictions precede either target GT convention. Reference annotation agreement is not complete-object ground truth. Original negative scribbles and target RGB are explicitly declared inputs. Other four benchmarks remain unchanged references.
\begin{center}\small\resizebox{\linewidth}{!}{\begin{tabular}{lrrrrr}
\toprule
Variant & IoU & AP & Precision & Recall & LUDVIG GT IoU \\\midrule
v10 & 86.51 & 98.02 & 95.70 & 90.10 & 86.68 \\
Static image & 89.06 & 95.04 & 96.13 & 92.51 & 89.13 \\
Static native & 88.84 & 96.58 & 96.21 & 92.20 & 88.96 \\
Legacy memory image & 88.40 & 98.66 & 94.97 & 93.10 & 88.39 \\
Legacy memory native & 88.20 & 98.35 & 95.10 & 92.73 & 88.27 \\
Semantic memory image & 88.81 & 97.00 & 91.08 & 96.96 & 88.75 \\
Semantic memory native & 88.52 & 96.29 & 91.17 & 96.58 & 88.56 \\
Constrained memory image & 91.68 & 99.00 & 94.26 & 97.13 & 91.59 \\
Constrained memory native (primary) & 91.48 & 98.78 & 94.35 & 96.79 & 91.47 \\
\bottomrule
\end{tabular}
}\end{center}
\begin{center}\small\begin{tabular}{lrrrr}\toprule
Task & Static & Legacy memory & Semantic memory & Constrained memory \\\midrule
fern & 83.92 & 84.06 & 83.20 & 83.25 \\
flower & 85.98 & 89.25 & 96.71 & 96.60 \\
fortress & 97.36 & 97.66 & 98.00 & 97.80 \\
horns\_center & 97.30 & 95.47 & 95.69 & 97.12 \\
horns\_left & 93.68 & 94.80 & 94.81 & 84.48 \\
leaves & 95.61 & 95.68 & 57.63 & 94.84 \\
orchids & 73.17 & 70.59 & 95.60 & 94.52 \\
trex & 83.71 & 78.09 & 86.54 & 83.21 \\
\bottomrule\end{tabular}\end{center}
constrained memory native vs static native: +2.6350 IoU points, +2.1944 AP points; 3 improved and 5 harmed tasks.
constrained memory native vs legacy memory native: +3.2772 IoU points, +0.4211 AP points; 5 improved and 3 harmed tasks.
constrained memory native vs semantic memory native: +2.9548 IoU points, +2.4814 AP points; 3 improved and 5 harmed tasks.
semantic memory native vs legacy memory native: +0.3224 IoU points, -2.0603 AP points; 6 improved and 2 harmed tasks.
constrained memory native vs constrained memory image: -0.2046 IoU points, -0.2275 AP points; 2 improved and 6 harmed tasks.
\begin{figure}[ht]\centering\includegraphics[width=.85\linewidth,height=.75\textheight,keepaspectratio]{../../output/tmm_response_20261018/study/nvos_reference_completion_spatial_v1/errors.png}
\caption{All eight N12 tasks. Green TP, red FP, blue FN; native-grid display resizing only. Scores use original GT resolution.}\end{figure}
}{}
\IfFileExists{tables/nvos_reference_delta.tex}{\section{Matched reference-intervention residuals}
N13 was fixed before N12 target scores. It reuses the exact static response b and paired original/completed-reference memory responses p and q. The primary response is clip(b+q-p,0,1); a fixed secondary uses the equivalent log-odds difference. Unchanged paired teacher values retain b exactly. The intervention may add or remove foreground and does not guarantee correctness. All branches undergo actual native-ray 3D fitting; 56 predictions precede target GT. Full N12 replacement is the matched control; earlier N11 is a separate candidate, not an additive causal comparison.
\begin{center}\small\resizebox{\linewidth}{!}{\begin{tabular}{lrrrrr}
\toprule
Variant & IoU & AP & Precision & Recall & LUDVIG GT IoU \\\midrule
v10 & 86.51 & 98.02 & 95.70 & 90.10 & 86.68 \\
Static native & 88.84 & 96.58 & 96.21 & 92.20 & 88.96 \\
Full replacement native & 91.48 & 98.78 & 94.35 & 96.79 & 91.47 \\
Probability delta image & 90.91 & 95.74 & 94.10 & 96.40 & 90.88 \\
Probability delta native (primary) & 90.75 & 97.04 & 94.18 & 96.12 & 90.78 \\
Log-odds delta image & 89.45 & 95.47 & 93.66 & 95.30 & 89.43 \\
Log-odds delta native & 89.44 & 96.74 & 93.77 & 95.17 & 89.48 \\
\bottomrule
\end{tabular}
}\end{center}
\begin{center}\small\begin{tabular}{lrrrr}\toprule
Task & Static & Replacement & Probability delta & Log-odds delta \\\midrule
fern & 83.92 & 83.25 & 83.14 & 81.07 \\
flower & 85.98 & 96.60 & 94.51 & 94.57 \\
fortress & 97.36 & 97.80 & 97.52 & 97.59 \\
horns\_center & 97.30 & 97.12 & 96.57 & 94.78 \\
horns\_left & 93.68 & 84.48 & 83.71 & 83.29 \\
leaves & 95.61 & 94.84 & 94.71 & 94.63 \\
orchids & 73.17 & 94.52 & 94.40 & 88.63 \\
trex & 83.71 & 83.21 & 81.41 & 80.98 \\
\bottomrule\end{tabular}\end{center}
probability native vs static native: +1.9045 IoU points, +0.4596 AP points; 3 improved and 5 harmed tasks.
probability native vs constrained memory native: -0.7305 IoU points, -1.7347 AP points; 0 improved and 8 harmed tasks.
probability native vs log odds native: +1.3053 IoU points, +0.3000 AP points; 6 improved and 2 harmed tasks.
probability native vs probability image: -0.1680 IoU points, +1.3063 AP points; 1 improved and 7 harmed tasks.
\begin{figure}[ht]\centering\includegraphics[width=.85\linewidth,height=.75\textheight,keepaspectratio]{../../output/tmm_response_20261018/study/nvos_reference_delta_spatial_v1/errors.png}
\caption{All eight N13 tasks. Green TP, red FP, blue FN; native-grid display resizing only. Scores use original GT resolution.}\end{figure}
}{}
\IfFileExists{tables/nvos_reference_admission.tex}{\section{Source-conditioned reference admission}
N14 is an adaptive development study designed after N12 analysis. A fixed feasibility policy retains semantic completion when it satisfies every original signed annotation, otherwise using constrained completion. A source-quality control instead chooses the larger SAM-predicted reference mask quality, with deterministic semantic-first ties. Both use only reference evidence, without scene names or target scores. Each selected target prediction is freshly replayed and hash-checked against N12, then actually fitted and re-rendered in 3D. All 56 predictions precede scoring. This is not an independent final test or a task-wise oracle. Other four benchmarks remain unchanged references.
\begin{center}\small\resizebox{\linewidth}{!}{\begin{tabular}{lrrrrr}
\toprule
Variant & IoU & AP & Precision & Recall & LUDVIG GT IoU \\\midrule
v10 & 86.51 & 98.02 & 95.70 & 90.10 & 86.68 \\
Static image & 89.06 & 95.04 & 96.13 & 92.51 & 89.13 \\
Static native & 88.84 & 96.58 & 96.21 & 92.20 & 88.96 \\
Feasibility image & 93.22 & 99.08 & 95.72 & 97.20 & 93.17 \\
Feasibility native (primary) & 92.94 & 98.86 & 95.78 & 96.84 & 92.99 \\
Source quality image & 93.50 & 99.13 & 96.09 & 97.12 & 93.43 \\
Source quality native (secondary) & 93.17 & 98.95 & 96.15 & 96.71 & 93.21 \\
\bottomrule
\end{tabular}
}\end{center}
\begin{center}\small\begin{tabular}{lrrrr}\toprule
Task & v10 & Static & Feasibility & Source quality \\\midrule
fern & 82.19 & 83.92 & 83.25 & 83.20 \\
flower & 86.53 & 85.98 & 96.71 & 96.71 \\
fortress & 97.23 & 97.36 & 98.00 & 98.00 \\
horns\_center & 95.73 & 97.30 & 97.12 & 95.69 \\
horns\_left & 92.70 & 93.68 & 94.81 & 94.81 \\
leaves & 94.08 & 95.61 & 94.84 & 94.84 \\
orchids & 68.14 & 73.17 & 95.60 & 95.60 \\
trex & 75.51 & 83.71 & 83.21 & 86.54 \\
\bottomrule\end{tabular}\end{center}
feasible native vs static native: +4.0999 IoU points, +2.2782 AP points; 4 improved and 4 harmed tasks.
quality native vs static native: +4.3305 IoU points, +2.3634 AP points; 5 improved and 3 harmed tasks.
quality native vs feasible native: +0.2306 IoU points, +0.0852 AP points; 1 improved and 2 harmed tasks.
feasible native vs feasible image: -0.2776 IoU points, -0.2173 AP points; 1 improved and 7 harmed tasks.
quality native vs quality image: -0.3241 IoU points, -0.1882 AP points; 1 improved and 7 harmed tasks.
\begin{figure}[ht]\centering\includegraphics[width=.85\linewidth,height=.75\textheight,keepaspectratio]{../../output/tmm_response_20261018/study/nvos_reference_admission_spatial_v1/errors.png}
\caption{All eight N14 tasks. Green TP, red FP, blue FN; native-grid display resizing only. Scores use original GT resolution.}\end{figure}
}{}
\IfFileExists{tables/nvos_multiview_completion.tex}{\section{Completed-reference multiview fusion}
N15 freezes N14 reference-quality selection and independently transports original and completed reference masks to every authorized source view with the same two-frame memory. Source-only fusion excludes target image evidence; the joint arm balances the average normalized source objective with one native target objective. All three columns use identical geometry, initial responses, observed elements and 120 Huber iterations. Source evidence is integrated at $512\times384$, with native $1008\times756$ target rays and rendering. There is no view-subset or fusion-weight search. All 48 predictions precede both GT scoring conventions. The original eight tasks are development data. Other four benchmarks remain unchanged references.
\begin{center}\small\resizebox{\linewidth}{!}{\begin{tabular}{lrrr}
\toprule
Variant & IoU & AP & LUDVIG GT IoU \\\midrule
v10 & 86.51 & 98.02 & 86.68 \\
N14 image & 93.50 & 99.13 & 93.43 \\
N14 native & 93.17 & 98.95 & 93.21 \\
Original source fusion & 88.01 & 98.26 & 88.15 \\
Completed source fusion & 92.75 & 98.81 & 92.90 \\
Completed joint fusion (primary) & 93.03 & 98.98 & 93.15 \\
\bottomrule
\end{tabular}
}\end{center}
\begin{center}\small\begin{tabular}{lrrrr}\toprule
Task & N14 & Original source & Complete source & Complete joint \\\midrule
fern & 83.20 & 83.72 & 82.92 & 83.68 \\
flower & 96.71 & 87.55 & 95.77 & 96.29 \\
fortress & 98.00 & 97.45 & 97.70 & 97.76 \\
horns\_center & 95.69 & 96.52 & 95.32 & 95.48 \\
horns\_left & 94.81 & 95.05 & 95.09 & 94.78 \\
leaves & 94.84 & 94.97 & 94.66 & 94.70 \\
orchids & 95.60 & 68.49 & 94.40 & 94.99 \\
trex & 86.54 & 80.30 & 86.12 & 86.57 \\
\bottomrule\end{tabular}\end{center}
completed source native vs legacy source native: +4.7400 IoU points, +0.5563 AP points; 5 improved and 3 harmed tasks.
completed joint native vs completed source native: +0.2841 IoU points, +0.1648 AP points; 7 improved and 1 harmed tasks.
completed joint native vs quality native: -0.1428 IoU points, +0.0346 AP points; 2 improved and 6 harmed tasks.
completed source native vs quality native: -0.4269 IoU points, -0.1302 AP points; 1 improved and 7 harmed tasks.
The matched source-only comparison isolates reference completion under the same transport and fusion, with zero target-image evidence in either field. Its gain remains +1.7152 to +6.5129 points when leaving out one whole scene. This is descriptive cohort sensitivity, not independent validation. The historical N01 comparison additionally changes negative-scribble and DINOv2 access and is not a matched control.
Against N14, joint fusion reduces mean completed-source Huber residual from 0.004626 to 0.000748. Mean precision changes by +0.5054 points and recall by -0.6629. Teacher agreement therefore does not by itself certify target accuracy. These teacher residuals are not source-view ground-truth scores.
The predeclared stopping criterion requires higher mean IoU without lower AP. No arm qualifies; this fusion branch is stopped and N14 is retained. This is a local negative result, not a universal performance bound.
\begin{figure}[ht]\centering\includegraphics[width=.85\linewidth,height=.75\textheight,keepaspectratio]{../../output/tmm_response_20261018/study/nvos_multiview_completion_spatial_v1/errors.png}
\caption{All eight N15 tasks: green TP, red FP, blue FN. Display resizing only; metrics use original GT resolution.}\end{figure}
}{}
\IfFileExists{sections/nvos_target_redecode.tex}{\section{Complete-query target re-decoding}
N16 replays N14 reference-quality completion and transport. The full-frame control uses that complete target response as one additional frozen SAM3 dense-logit and box prompt. The local primary applies the same operation to the response extent with 25 percent context on each side. RGB and prior share the same crop, and all outside probabilities are preserved before identical native 3D fitting. No extra annotations, candidate selection, recursive decoding or crop-size search is used. All 56 predictions precede both GT tracks. This is development on the opened eight-task cohort, with the other four benchmarks unchanged.
\begin{center}\small\begin{tabular}{lrrr}
\toprule
Variant & IoU & AP & LUDVIG GT IoU \\\midrule
v10 & 86.51 & 98.02 & 86.68 \\
N14 image & 93.50 & 99.13 & 93.43 \\
N14 native & 93.17 & 98.95 & 93.21 \\
Full-frame image & 93.36 & 97.14 & 93.34 \\
Full-frame native & 93.10 & 98.26 & 93.18 \\
Local image & 93.57 & 96.92 & 93.52 \\
Local native (primary) & 93.28 & 98.47 & 93.34 \\
\bottomrule
\end{tabular}
\end{center}
\begin{center}\small\begin{tabular}{lrrr}\toprule
Task & N14 native & Full native & Local native \\\midrule
fern & 83.20 & 83.68 & 83.68 \\
flower & 96.71 & 96.00 & 96.02 \\
fortress & 98.00 & 97.54 & 97.62 \\
horns\_center & 95.69 & 97.42 & 97.42 \\
horns\_left & 94.81 & 93.82 & 94.01 \\
leaves & 94.84 & 95.63 & 95.54 \\
orchids & 95.60 & 95.73 & 95.88 \\
trex & 86.54 & 84.98 & 86.08 \\
\bottomrule\end{tabular}\end{center}
The fixed window covers the full image on fern, horns\_center; local and full-frame image predictions there have identical hashes. These tasks do not test an image-scale change.
full native vs quality native: -0.0732 IoU points, -0.6824 AP points; 4 improved and 4 harmed tasks.
crop native vs quality native: +0.1075 IoU points, -0.4787 AP points; 4 improved and 4 harmed tasks.
crop native vs full native: +0.1806 IoU points, +0.2037 AP points; 5 improved and 1 harmed tasks.
full image vs quality image: -0.1367 IoU points, -1.9905 AP points; 4 improved and 4 harmed tasks.
crop image vs quality image: +0.0680 IoU points, -2.2103 AP points; 4 improved and 4 harmed tasks.
The fixed stopping criterion requires higher mean IoU without lower AP on both GT tracks. No arm qualifies; this target-refinement branch is stopped and N14 is retained. This does not establish a universal performance ceiling.
\begin{figure}[ht]\centering\includegraphics[width=.85\linewidth,height=.75\textheight,keepaspectratio]{../../output/tmm_response_20261018/study/nvos_target_redecode_spatial_v1/errors.png}
\caption{All eight N16 tasks: green TP, red FP, blue FN. Display resizing only; metrics use original GT resolution.}\end{figure}
}{}
\IfFileExists{sections/nvos_score_precision.tex}{\section{Response-score numerical audit}
A24 exactly replays all 24 N14/N16 image probability arrays and decoder decisions, while observing raw logits. All eight tasks are sealed before scoring on both GT conventions. The original float32 probability readout is reproduced, not replaced. Float64 evaluates the same sigmoid and bilinear operation at higher precision. The stable-mixture score independently interpolates positive and complementary probabilities, then takes their log ratio. It preserves fine upper-tail differences without changing to logit interpolation. Native-grid AP uses separately nearest-resized GT and remains diagnostic only.
\begin{center}\small\begin{tabular}{lrrr}\toprule
Image response & Original AP & Float64 AP & Stable mixture AP \\\midrule
N14 transport & 99.1333 & 99.1352 & 99.1352 \\
N16 full frame & 97.1428 & 97.1427 & 97.1427 \\
N16 local & 96.9230 & 96.9228 & 96.9228 \\
\bottomrule\end{tabular}
\end{center}
These are numerical diagnostics of image responses, not new 3D method scores. Actual 3D predictions are not recomputed; the N16 native AP loss and predeclared stopping decision remain unchanged.
quality image: stable minus original AP +0.0019 points; relative to matched N14 image AP, original +0.0000, stable +0.0000 points. Stable-readout hard-mask changes: 0 pixels across all tasks.
full image: stable minus original AP -0.0001 points; relative to matched N14 image AP, original -1.9905, stable -1.9925 points. Stable-readout hard-mask changes: 0 pixels across all tasks.
crop image: stable minus original AP -0.0002 points; relative to matched N14 image AP, original -2.2103, stable -2.2124 points. Stable-readout hard-mask changes: 0 pixels across all tasks.
\begin{center}\small\begin{tabular}{lrrr}\toprule
Task & N14 AP correction & Full AP correction & Local AP correction \\\midrule
fern & +0.0151 & +0.0000 & +0.0000 \\
flower & +0.0000 & -0.0000 & -0.0008 \\
fortress & +0.0000 & +0.0000 & +0.0000 \\
horns\_center & +0.0000 & +0.0000 & +0.0000 \\
horns\_left & +0.0000 & -0.0007 & -0.0008 \\
leaves & +0.0000 & +0.0000 & +0.0000 \\
orchids & +0.0000 & -0.0000 & -0.0000 \\
trex & +0.0000 & -0.0003 & -0.0001 \\
\bottomrule\end{tabular}\end{center}
}{}
\clearpage
\section{Complete NVOS experiment catalog}
The main paper uses a compact editorial view. This table preserves every
registered NVOS comparison, including failed or superseded variants and the
independent selected-set readout. P/R/O/M denote published, reproduced, our
established implementation and measured extensions. Teacher, prompt, target-RGB
and evaluation contracts differ as stated in the individual study sections;
these rows are not an input-identical ranking.
\begin{center}\small\begin{tabular}{llr}
\toprule
Method & Origin & IoU \\
\midrule
NVOS \cite{ludvig} & P & 70.10 \\
SA3D-TRF \cite{ludvig} & P & 90.30 \\
SA3D-GS \cite{ludvig} & P & 92.20 \\
SAGA \cite{ludvig} & P & 92.60 \\
OmniSeg3D \cite{ludvig} & P & 91.70 \\
LUDVIG-DINOv2 \cite{ludvig} & P & 92.40 \\
LUDVIG-SAM \cite{ludvig} & P & 91.30 \\
LUDVIG-SAM2 \cite{ludvig} & P & 91.30 \\
VCAR \cite{vcar} & P & 93.50 \\
LEGO \cite{lego} & P & 94.20 \\
WildSeg3D \cite{wildseg3d} & P & 94.10 \\
Trace3D \cite{trace3d} & P & 92.50 \\
LUDVIG-SAM & R & 91.26 \\
Ours (v10) & O & 86.51 \\
Ours + coverage (N01) & M & 87.78 \\
Ours + coverage + source readout (N02) & M & 87.58 \\
Ours + local edit credit (N03) & M & 87.55 \\
Ours + image response + 3D (N04) & M & 88.32 \\
Ours + partial image response + 3D (N05) & M & 85.38 \\
Ours + unboxed image response + 3D (N06) & M & 84.15 \\
Ours + reference memory + 3D (N07) & M & 87.89 \\
Ours + native memory writeback (N08) & M & 88.11 \\
Ours + static native writeback (N08 secondary) & M & 88.84 \\
Ours + signed reference affinity (N09) & M & 81.58 \\
Ours + direct signed affinity (N10) & M & 77.73 \\
Ours + instance-constrained extent (N11) & M & 89.97 \\
Ours + signed reference completion (N12) & M & 91.48 \\
Ours + paired reference delta (N13) & M & 90.75 \\
Ours + source feasibility admission (N14) & M & 92.94 \\
Ours + source quality admission (N14 secondary) & M & 93.17 \\
Ours + completed multiview fusion (N15) & M & 93.03 \\
Ours + query-local target decoding (N16) & M & 93.28 \\
Ours + full-frame target decoding (N16 control) & M & 93.10 \\
Ours (v10, selected set) & O & 85.22 \\
\bottomrule
\end{tabular}
\end{center}
\section{Static image-response control: implementation}
\paragraph{Image-conditioned response refinement.}
For an observed view $t$, let $z_t^0=A_t s^0$ be the rendered v10 response.
We use its thresholded extent as a box and its bounded log-odds as a dense
prompt for a frozen image decoder $D$:
\begin{equation}
 p_t=D(I_t,\operatorname{box}[z_t^0\geq0.5],
       \operatorname{resize}[\operatorname{logit}(\operatorname{clip}(z_t^0))]).
\end{equation}
Clipping uses $[10^{-4},1-10^{-4}]$; the dense prompt is resized to the
decoder's actual $288\times288$ grid. The box and dense mask are fallible
predictions, not extra annotations. Single-mask decoding avoids target-label
candidate selection. Paired reconstructed/observed RGB controls isolate image
access under this fixed interface. Sparse-point and unboxed controls test
whether the inherited response preserves identity while restricting extent.


This is the earlier static control, not an additional stage of the
reference-completion pipeline. Coarse fits pool image evidence on a
$512\times384$ grid; the native control uses $1008\times756$ rays.

\begin{thebibliography}{9}
\bibitem{vala}
S. Wang, K. Li, S. Liang, E. Alegret, J. Ma, N. Navab, and S. Gasperini.
\emph{Visibility-Aware Language Aggregation for Open-Vocabulary Segmentation
in 3D Gaussian Splatting}. arXiv:2509.05515v1, 2025.
\url{https://arxiv.org/abs/2509.05515v1}.
Comparison values are transcribed from Tables 1 and 2; their cited
baselines retain the reporting paper's provenance.
\bibitem{ludvig}
J. Marrie, R. Menegaux, M. Arbel, D. Larlus, and J. Mairal.
\emph{LUDVIG: Learning-Free Uplifting of 2D Visual Features to Gaussian
Splatting Scenes}. ICCV, 2025; arXiv:2410.14462v5.
\url{https://arxiv.org/abs/2410.14462v5}.
NVOS comparison values are transcribed from Tables 1 and D.
\bibitem{occam}
J. Cheng, J.-N. Zaech, L. Van Gool, and D. P. Paudel.
\emph{Occam's LGS: An Efficient Approach for Language Gaussian Splatting}.
arXiv:2412.01807v2, 2025.
\url{https://arxiv.org/abs/2412.01807v2}.
\bibitem{evisplat}
S. Moon, K. Shimomura, J. Park, W. Choi, S. Lee, T. Yamashita, and S. Im.
\emph{EviSplat: Preserving Multi-View Evidence in 3D Gaussian Splatting for
Open-Vocabulary Segmentation}. arXiv:2609.34853v1, 2026.
\url{https://arxiv.org/abs/2609.34853v1}.
\bibitem{sam3}
N. Carion et al. \emph{SAM 3: Segment Anything with Concepts}.
arXiv:2511.16719, 2025. \url{https://arxiv.org/abs/2511.16719}.
\bibitem{siglip2}
M. Tschannen et al. \emph{SigLIP 2: Multilingual Vision-Language Encoders with
Improved Semantic Understanding, Localization, and Dense Features}.
arXiv:2502.14786, 2025. \url{https://arxiv.org/abs/2502.14786}.
\bibitem{dinov2}
M. Oquab et al. \emph{DINOv2: Learning Robust Visual Features without Supervision}.
arXiv:2304.07193, 2023. \url{https://arxiv.org/abs/2304.07193}.
\bibitem{maskadapter}
Y. Li, T. Cheng, B. Feng, W. Liu, and X. Wang.
\emph{Mask-Adapter: The Devil is in the Masks for Open-Vocabulary Segmentation}.
CVPR, 2025. \url{https://github.com/hustvl/MaskAdapter}.
\bibitem{vcar}
K. Cao, D. Wang, H. Zhu, H. Huang, X. Wang, Z. Shi, and G. Yang.
\emph{VCAR: Training-Free 3DGS Segmentation via View Completeness and Axis-Aware Boundary Refinement}.
arXiv:2608.30870v1, 2026. \url{https://arxiv.org/abs/2608.30870v1}.
\bibitem{lego}
Y. Peng, H. Wang, Y. Liu, Y. Lu, Z. Dong, and B. Yang.
\emph{LEGO: Leveled Language Gaussian Splatting}.
ECCV, 2026. \url{https://arxiv.org/abs/2608.10057v1}.
\bibitem{wildseg3d}
Y. Guo, J. Hu, Y. Qu, and L. Cao.
\emph{WildSeg3D: Segment Any 3D Objects in the Wild from 2D Images}.
ICCV, 2025. \url{https://arxiv.org/abs/2503.08407}.
\bibitem{trace3d}
H. Shen, J. Ni, Y. Chen, W. Li, M. Pei, and S. Huang.
\emph{Trace3D: Consistent Segmentation Lifting via Gaussian Instance Tracing}.
ICCV, 2025. \url{https://arxiv.org/abs/2508.03227v1}.
\end{thebibliography}

\end{document}
